\section{Results}
\label{sec:results}

\subsection{RQ1: weapon-space composition}

\Cref{tab:formal-candidates} reports the four paired candidates. \routeB\ reduces
asset-seam error from 0.2438--0.2967 to below 0.0008 for every pair. It also
improves the multi-view score for every pair, despite a modest reduction in the
texture-only score. Candidate 04 is selected within \routeB.

\begin{table}[t]
  \centering
  \caption{Formal paired A/B results. Lower asset-seam error is better; higher
  texture and multi-view scores are better.}
  \label{tab:formal-candidates}
  \small
  \begin{tabular}{llrrrr}
    \toprule
    Route & Candidate & Asset seam & Multi-view & Texture & Selected\\
    \midrule
    A & 01 & 0.286821 & 0.752696 & 0.761178 & No\\
    A & 02 & 0.252571 & 0.792474 & 0.779769 & Yes\\
    A & 03 & 0.243767 & 0.767358 & 0.770474 & No\\
    A & 04 & 0.296737 & 0.763398 & 0.776778 & No\\
    \midrule
    B & 01 & 0.000473 & 0.819614 & 0.765960 & No\\
    B & 02 & 0.000635 & 0.818653 & 0.749159 & No\\
    B & 03 & 0.000617 & 0.833842 & 0.755199 & No\\
    B & 04 & 0.000796 & 0.844312 & 0.751902 & Yes\\
    \bottomrule
  \end{tabular}
\end{table}

The paired summary in \cref{tab:paired-summary} shows a mean asset-seam
improvement of 0.26934 and a mean multi-view improvement of 0.06012. The win rate
is 100\% on both measures. The bootstrap intervals remain above zero for the four
observed pairs.

\begin{table}[t]
  \centering
  \caption{Paired Route B minus Route A summary ($n=4$).}
  \label{tab:paired-summary}
  \small
  \begin{tabular}{lrrrrr}
    \toprule
    Metric & Mean & Median & Std. & Win rate & Bootstrap 95\% CI\\
    \midrule
    Asset-seam improvement & 0.269344 & 0.269142 & 0.022281 & 100\% &
    [0.247543, 0.291145]\\
    Multi-view improvement & 0.060124 & 0.066701 & 0.020440 & 100\% &
    [0.036363, 0.077307]\\
    \bottomrule
  \end{tabular}
\end{table}

\begin{figure}[t]
  \centering
  \begin{subfigure}[t]{0.49\textwidth}
    \includegraphics[width=\textwidth]{formal_route_a_candidate04.png}
    \caption{Route A candidate 04}
  \end{subfigure}
  \hfill
  \begin{subfigure}[t]{0.49\textwidth}
    \includegraphics[width=\textwidth]{formal_route_b_candidate04.png}
    \caption{Route B candidate 04}
  \end{subfigure}
  \caption{Mapped left, right, and top views for the same candidate index in the
  formal experiment. Route A behaves as an all-over pattern; Route B reconstructs
  a continuous composition after UV fragmentation.}
  \label{fig:formal-ab}
\end{figure}

The formal seam-correction operation reduces the mean \routeB\ seam error by
99.62\%. Before correction, the four values range from 0.14986 to 0.17639; after
correction, all are below 0.0008 and pass the 0.01 threshold. This establishes that
the low final seam error is not merely inherited from the source artwork.

These results answer RQ1 positively for the tested AK-47 candidates. The conclusion
is bounded to this asset, theme family, renderer, and four-pair sample.

\subsection{RQ2: failures exposed by mapped evaluation}

The live dragon case provides direct examples of why source images are
insufficient. One attempted artwork passed technical image checks but was rejected
because its subjects formed isolated clusters around large dark regions. Such a
composition is fragile under cropping and UV fragmentation. After a local retry,
all four candidate sources were mapped before selection.

\begin{figure}[t]
  \centering
  \begin{subfigure}[t]{0.98\textwidth}
    \includegraphics[width=\textwidth]{selected_master_artwork.png}
    \caption{Human-selected landscape source artwork}
  \end{subfigure}
  \vspace{2mm}
  \begin{subfigure}[t]{0.98\textwidth}
    \includegraphics[width=\textwidth]{dragon_route_b_multiview.png}
    \caption{Formal Route B left, right, and top reconstruction}
  \end{subfigure}
  \caption{The selected source and its mapped result are related but not
  interchangeable evidence. The narrow top view and UV fragmentation remove or
  compress details that appear clear in the source.}
  \label{fig:dragon-selected}
\end{figure}

For the selected source, automatic review reports source fulfillment 0.9386, left
readability 0.6016, right readability 0.6016, and top readability 0.3993. Only
eight of sixteen planned elements exceed the visibility threshold. The top view
is the weakest because of its narrow projected area, while the source score is
much higher. This difference is the central observation for RQ2: mapping reveals
loss of readable elements, view imbalance, component-specific cropping, and
seam behavior that cannot be inferred reliably from the rectangular artwork.

The route comparison in \cref{fig:dragon-routes} illustrates a second failure
mode. The generated \routeA\ image has a high multi-view score of 0.8952, but its
asset-seam error is 0.1822 and it fails the deployment constraint. A strong
aggregate visual score is therefore not evidence of a technically valid texture.

\begin{figure}[t]
  \centering
  \begin{subfigure}[t]{0.49\textwidth}
    \includegraphics[width=\textwidth]{dragon_route_a_multiview.png}
    \caption{Route A: strong visible activity, invalid asset seam}
  \end{subfigure}
  \hfill
  \begin{subfigure}[t]{0.49\textwidth}
    \includegraphics[width=\textwidth]{dragon_route_b_multiview.png}
    \caption{Route B: valid seam and continuous mapped composition}
  \end{subfigure}
  \caption{The final live case separates visual activity from deployment
  feasibility.}
  \label{fig:dragon-routes}
\end{figure}

\subsection{RQ3: constrained selection and rollback}

\Cref{fig:edge-sweep} shows the expected trade-off in the independent edge-width
study. Increasing the safety width reduces seam error but gradually removes
multi-view detail. Widths 8, 12, and 16 are feasible and Pareto-optimal. Width 24
has a lower seam error but falls below 95\% retention.

\begin{figure}[t]
  \centering
  \begin{tikzpicture}
    \begin{axis}[
      width=0.92\textwidth,
      height=62mm,
      xlabel={Edge-safety width (pixels)},
      ylabel={Score},
      xmin=0,xmax=24,
      ymin=0,ymax=1.02,
      grid=both,
      legend style={at={(0.5,-0.22)},anchor=north,legend columns=3},
      mark size=2.3pt
    ]
      \addplot+[thick,mark=*,skinsmithblue]
        table[x=edge_safe_pixels,y=asset_seam_error,col sep=comma]
        {data/table_edge_safety_sweep.csv};
      \addlegendentry{Asset-seam error}
      \addplot+[thick,mark=square*,skinsmithteal]
        table[x=edge_safe_pixels,y=multi_view_score,col sep=comma]
        {data/table_edge_safety_sweep.csv};
      \addlegendentry{Multi-view score}
      \addplot+[thick,mark=triangle*,orange]
        table[x=edge_safe_pixels,y=multi_view_retention,col sep=comma]
        {data/table_edge_safety_sweep.csv};
      \addlegendentry{Multi-view retention}
      \addplot[red,dashed,domain=0:24] {0.01};
      \addplot[gray,dashed,domain=0:24] {0.95};
    \end{axis}
  \end{tikzpicture}
  \caption{Independent edge-width sweep. Width 8 is the selected feasible
  operating point. The two dashed lines show the seam and retention thresholds.}
  \label{fig:edge-sweep}
\end{figure}

The policy ablation is decisive. Equal weighted ranking selects width 4 because it
retains a multi-view score of 0.70324, but its seam error is 0.01660 and therefore
invalid. Constraint-first/Pareto selection chooses width 8, with seam error
0.00552 and multi-view score 0.69697. The weighted approach hides a hard violation
behind a small visual gain.

All four formal \routeC\ cases execute two local corrections. The selected
attempts change 2.33--2.93\% of atlas pixels and change zero pixels outside the
target component plus halo. Their best locked-score changes are -0.00110,
-0.00110, -0.00082, and -0.00088, all below the required +0.01. Consequently,
every final C result equals B. In the live dragon case, the best correction changes
the score by -0.02196 and is also rolled back.

RQ3 is therefore supported in two ways: feasibility-first selection rejects a
candidate preferred by a soft score, and the acceptance gate prevents a lower
quality correction from replacing its baseline. The result does not show that
refinement improves quality; it shows that attempted refinement is bounded.

\subsection{RQ4: three-checkpoint collaboration}

The live Agent completes the intended interaction:

\begin{enumerate}
  \item the theme is expanded into dragons, clouds, waves, treasure, flaming
  pearls, mountains, lightning, mist, jade, and architectural fragments;
  \item four directions are presented, and \emph{Treasured Relic} is selected;
  \item four artworks are generated and mapped, and \texttt{artwork\_04} is
  selected from its source-plus-three-view card.
\end{enumerate}

The run uses six image calls, one local retry, and one refinement round. It ends in
the \texttt{completed} state and exports a 2048 by 2048 RGB PNG and TGA with
identical decoded pixels. The final \routeB\ result has asset-seam error 0.000778,
multi-view score 0.816620, and total score 0.829123. The source hash and final TGA
hash are preserved.

This case answers RQ4 at the level of functional workflow: a user can begin from a
compact theme, make three distinct decisions, and receive a formally processed
texture without authoring a production prompt. It does not constitute a usability
study; no completion-time or user-satisfaction comparison was conducted.

\subsection{RQ5: adapter-based portability}

The M4A4 adapter executes geometry mapping, continuous composition, UV baking,
seam measurement, rendering, and evaluation with no new image-generation call.
\Cref{fig:m4-transfer} confirms that the same conceptual pipeline reconstructs a
mapped texture on different geometry.

\begin{figure}[t]
  \centering
  \includegraphics[width=\textwidth]{m4a4_transfer_concept.png}
  \caption{M4A4 adapter transfer. The core method is reused, but the denser UV
  topology requires different seam calibration.}
  \label{fig:m4-transfer}
\end{figure}

Transfer is only partial. The best joint candidate reaches seam error 0.00994 but
retains 94.02\% of the raw multi-view score, below the AK-47 threshold of 95\%.
The M4A4 contains 1,747 UV islands and 11,341 seam pairs, compared with 928 islands
and 4,999 pairs for the AK-47. More aggressive seam treatment improves continuity
but removes too much visible detail.

RQ5 is therefore answered by a boundary rather than a universal success:
planning, composition, bake, rendering, metrics, and logging transfer; mesh, UV,
axes, camera, material, and correction calibration remain adapter-specific.
